\documentclass[12pt]{article}
\usepackage[T1]{fontenc}
\usepackage[utf8]{inputenc}
\usepackage{lmodern}
\usepackage{textcomp}
\usepackage{enumitem}
\usepackage{geometry}
\geometry{margin=1in}
\begin{document}
\begin{center}
Answer Key - Psychology - Graduate Level Assessment 2 \
\small (Answers are ordered exactly as in the question paper.)
\end{center}

\section*{Multiple Choice}
\begin{enumerate}[label=\arabic*.]
\item \textbf{Answer:} A
\\ \textit{Options:} A) association.; B) generalization.; C) recognition.; D) classification.; E) interpretation.
\\ \textit{Note:} Antonia's mistake is due to the process of association, as she incorrectly identifies the rabbit as a cat based on her previous experiences and learned associations with animals, leading to a misclassification.
\item \textbf{Answer:} A
\\ \textit{Options:} A) less than 5 seconds.; B) 30 to 45 seconds.; C) only while the sensory input is present.; D) 10 to 15 seconds.; E) exactly 5 minutes.
\\ \textit{Note:} Sensory memory retains images and sounds for a brief duration, typically less than 5 seconds, allowing individuals to process and respond to sensory stimuli before they fade away.
\item \textbf{Answer:} A
\\ \textit{Options:} A) Neurotransmitters diffuse away and are exhaled; B) Neurotransmitters are converted into cellular energy; C) Neurotransmitters replicate to enhance the initial signal; D) Neurotransmitters are destroyed; E) Neurotransmitters are released into the bloodstream
\\ \textit{Note:} Neurotransmitters are released into the synaptic cleft, and after transmitting their signal, they diffuse away to prevent continuous stimulation of the postsynaptic neuron, which is crucial for the proper functioning of neural communication.
\item \textbf{Answer:} A
\\ \textit{Options:} A) solely by their physical needs; B) striving for self-actualization and personal growth; C) primary (physical) and secondary (emotional) needs; D) a constant state of dissatisfaction; E) desire for power and dominance over others
\\ \textit{Note:} The correct answer reflects a misunderstanding of Maslow's hierarchy, which posits that individuals are motivated by a range of needs, including physiological, safety, love and belonging, esteem, and self-actualization, rather than solely by their physical needs.
\item \textbf{Answer:} A
\\ \textit{Options:} A) Socioeconomic status and personal beliefs; B) Childhood upbringing and current employment status; C) Diet and exercise habits; D) Personality type and coping mechanisms; E) Intelligence quotient (IQ) and gender
\\ \textit{Note:} Socioeconomic status and personal beliefs are critical factors that influence an individual's resilience to conflict and their mental health outcomes, as they shape coping mechanisms, access to resources, and the overall perception of stressors.
\end{enumerate}

\section*{True / False}
\begin{enumerate}[label=\arabic*.]
\setcounter{enumi}{5}
\item \textbf{Answer:} False
\\ \textit{Options:} A) True; B) False
\\ \textit{Note:} People with Wernicke's aphasia typically produce fluent but nonsensical speech and exhibit significant difficulties in comprehension, rather than demonstrating nonfluent and relevant speech.
\item \textbf{Answer:} False
\\ \textit{Options:} A) True; B) False
\\ \textit{Note:} Whole-interval recording measures the occurrence of a behavior during specified intervals, indicating whether the behavior was present for the entire duration of each interval, rather than measuring the total time the behavior occurs.
\item \textbf{Answer:} False
\\ \textit{Options:} A) True; B) False
\\ \textit{Note:} Operant extinction is not always preferred because it can lead to an initial increase in the unwanted behavior, known as an extinction burst, which may exacerbate the issue before the behavior decreases.
\item \textbf{Answer:} True
\\ \textit{Options:} A) True; B) False
\\ \textit{Note:} Meditation promotes relaxation and a state of calm that leads to increased alpha and theta brain wave activity, while simultaneously reducing metabolic rates and heart rates due to its effects on the autonomic nervous system.
\item \textbf{Answer:} True
\\ \textit{Options:} A) True; B) False
\\ \textit{Note:} The statement is true because Hofstede's cultural dimensions indicate that the United States prioritizes personal autonomy, self-expression, and individual rights over collective goals, reflecting a strong individualistic culture.
\end{enumerate}

\section*{Long Form Response}
\begin{enumerate}[label=\arabic*.]
\setcounter{enumi}{10}
\item \textbf{Answer:} Attribution theory, which seeks to explain how individuals interpret the causes of behavior, is inherently influenced by stereotyping, as it relies on generalized beliefs about social groups. Stereotyping simplifies complex social information, allowing for quicker judgments about others based on perceived group characteristics. This reliance on stereotypes can lead to biased attributions, where individuals may attribute behaviors to inherent traits of a group rather than considering situational factors. For instance, if a person from a stereotyped group exhibits a negative behavior, observers might attribute this to the person's character rather than external circumstances. Thus, attribution theory operates within a framework shaped by stereotypes, illustrating how these generalizations can distort our understanding of individual actions and motivations.
\\ \textit{Note:} correct because it effectively illustrates how stereotyping shapes attribution processes by promoting biased interpretations of behavior based on generalized beliefs about social groups, leading to misattributions that overlook situational factors.
\item \textbf{Answer:} Standard scores and percentiles are critical tools in psychological testing that facilitate the comparison of individuals across different assessments and within the same test. Standard scores, such as z-scores, offer a way to understand an individual's performance relative to a normative population, indicating how many standard deviations a score is from the mean. Percentiles, on the other hand, rank a score relative to others, showing the percentage of scores that fall below a particular point. Together, these metrics allow psychologists to make meaningful interpretations of test results, enabling comparisons that account for variations in test difficulty and population characteristics, thus enhancing the overall assessment process.
\\ \textit{Note:} correct because it accurately describes how standard scores and percentiles provide essential frameworks for interpreting individual performance relative to normative data and for comparing scores across different assessments, thereby enhancing the understanding of psychological test results.
\end{enumerate}

\vfill
\noindent\textit{End of Answer Key}
\end{document}